% !TeX root = ../../main.tex
\makeatletter
\renewcommand{\compareexrow}[2]{%
  \par\vspace{0.24em}%
  \noindent
  \begin{tabular}[t]{@{}p{\compareexwd}@{\hspace{\compareexgap}}p{\compareexwd}@{}}#1 & #2 \\ \end{tabular}%
  \par\vspace{0.1em}%
}
\renewcommand{\dialogueblankprac}[1]{%
  \par\vspace{0.1em}%
  \noindent
  \ifshowanswer
    \fitblank@withspaces{\underline{\makebox[\linewidth][l]{\hspace{0.5em}{\color{blue}\small #1}\hspace{0.5em}}}}%
  \else
    \fitblank@withspaces{\underline{\makebox[\linewidth][l]{\hspace{0.5em}{\small\phantom{#1}}\hspace{0.5em}}}}%
  \fi
  \par\vspace{0.12em}%
}
\makeatother
\subsection{最上級のつくり方}
    最上級のつくり方は以下の\fitblankbf[]{3パターン}である。

    \begin{definitionbox}{\textbf{最上級のつくり方}}
        \begin{enumerate}[nosep, topsep=0pt, itemsep=0.4em]
            \item 基本は語句の語尾に\fitblankbf[]{st}、\fitblankbf[]{est}をつける
            \compareexrow{\comparecellshort{tall}{\fitblankbf[]{tallest}}}{\comparecellshort{easy}{\fitblankbf[]{easiest}}}
            \compareexrow{\comparecellshort{nice}{\fitblankbf[]{nicest}}}{\comparecellshort{hot}{\fitblankbf[]{hottest}}}
            \item\fitblankbf[]{長い単語}(母音が3つ以上)のときは単語の前に\fitblankbf[]{most}をつける
            \par\vspace{0.5em}
            \noindent ※\uwave{前に most をつけると最上級変化をしていることと同義なので st や est はつけない}
            \compareexrow{\comparecelllong{careful}{\fitblankbf[]{most careful}}}{\comparecelllong{beautiful}{\fitblankbf[]{most beautiful}}}
            \compareexrow{\comparecelllong{delicious}{\fitblankbf[]{most delicious}}}{\comparecelllong{interesting}{\fitblankbf[]{most interesting}}}
            \par\vspace{0.5em}
            \item\fitblankbf[]{不規則}に変化するもの
            \compareexrow{\comparecellja{good, well}{(\fitblankbf[]{良い}、\fitblankbf[]{上手な～})}{\fitblankbf[]{best}}}{\comparecellja{many, much}{(\fitblankbf[]{多い}、たくさんの～)}{\fitblankbf[]{most}}}
            \compareexrow{\comparecellja{bad, ill}{(\fitblankbf[]{悪い}、すぐれない～)}{\fitblankbf[]{worst}}}{\comparecellja{little}{(\fitblankbf[]{小さい}、\fitblankbf[]{少し})}{\fitblankbf[]{least}}}
        \end{enumerate}
    \end{definitionbox}

    \begin{techniquebox}{\textbf{比較の不規則変化表}}
        \begin{center}
        \renewcommand{\arraystretch}{1.12}
        \begin{tabular}{@{}c@{\hspace{3em}}c@{\hspace{3em}}c@{}}
            \hline
            \textbf{原級} & \textbf{比較級} & \textbf{最上級} \\
            \hline
            good, well &\fitblankbf[]{better} &\fitblankbf[]{best} \\
            many, much &\fitblankbf[]{more} & most \\
            bad, ill & worse &\fitblankbf[]{worst} \\
            little &\fitblankbf[]{less} &\fitblankbf[]{least} \\
            \hline
        \end{tabular}
        \end{center}
    \end{techniquebox}
    \vfill

    \subsection{練習問題}
    \pracrow{\praccell{1}{happy}{happiest}}{\praccell{2}{big}{biggest}}{\praccell{3}{useful}{most useful}}
    \vspace{2.5em}
\newpage
